\section*{第 \ref{ch:g1:hot-and-cold} 章——冷与热}

\begin{solution}{exo:g1:hot-and-cold:1}
依次是冰块、卧室里的空气、洗澡水，然后
是刚出锅的汤。
\end{solution}

\begin{solution}{exo:g1:hot-and-cold:2}
与热巧克力相比，自来水是\emph{冷}的（或更凉），
但与果汁相比，它是\emph{温热}的（或更温暖）。
\end{solution}

\begin{solution}{exo:g1:hot-and-cold:3}
例如：太阳、火焰、暖气片、你自己的身体
（或搓手）。太阳温暖着整个公园。
\end{solution}

\begin{solution}{exo:g1:hot-and-cold:4}
手掌贴着脸颊时觉得暖和：搓手使它们变暖了。
\end{solution}

\begin{solution}{exo:g1:hot-and-cold:5}
阳光下的冰块先熔化：太阳使它变暖，而热量正是使冰熔化的原因。
阴影里的那块得到的热量较少，因此能维持得更久。
\end{solution}

\begin{solution}{exo:g1:hot-and-cold:6}
积水就是水——它与原来那块冰块是同一种水。
坚硬的冰并没有消失，房间里的热量把它
熔化成了液态水。
\end{solution}

\begin{solution}{exo:g1:hot-and-cold:7}
在同一片阴影里待了整个上午后，滑梯和长椅一样冷。
金属只是用手\emph{摸起来}更冷——
正如金属勺和木勺的实验所示，触觉并不是可靠的判断。
\end{solution}

\begin{solution}{exo:g1:hot-and-cold:8}
用\emph{同一}只手，先后立即摸一只碗、再摸另一只碗，
并请一位朋友在听不到你回答的情况下来检查。
之所以必须用同一只手，是因为“三碗”小实验表明，
两只手对同一碗水的感受可能不一样。
\end{solution}

\begin{solution}{exo:g1:hot-and-cold:9}
莉亚说得对。外套不会产生热量，它只是阻止热量
穿过。雪人身上没有需要留住的热量——外套的作用是
\emph{把白天的暖空气挡在外面}，因此裹着外套的雪人
熔化得更慢，存在的时间也更长。
\end{solution}
